% Complete original-source quantitative return, 14 September 2026.
% Providers: GDR, GPR, QLG, EIQ, WGP and LER; their proof bodies accompany this source.
\section{The original Gamma Hankel return with a finite moving-shift error}
\label{sec:QGT}

\paragraph{Original objects and exact square pushforward.}
Retain the original packet integers and the original density:
\[
 k=4l+1,\quad e=1+k(m-1),\quad q=e(k+1)^2=2n,
 \quad l,m\ge1,\quad
 \sigma(y)=\frac{|\Gamma(1/4+iy/2)|^2}{2\pi}.
 \tag{QGT1}
\]
In particular $n\ge18$ and $n\ge18l$: indeed
$n/l\ge8l+8+2/l\ge18$, the last expression increasing for $l\ge1$.
For every real $a>-1/2$ define the actual even moments and their
Hankel determinant by
\[
 \mu_a=\int_{\mathbb R}|y|^{2a}\sigma(y)\,dy,
 \qquad H_s^{(a)}=\det[\mu_{a+i+j}]_{0\le i,j<s},\quad s\ge1.
 \tag{QGT2}
\]
For the integer values used in the programme this is precisely
$\mu_a=m_{2a}$, with $\mu_0=\sqrt{2\pi}$ retained. The extension to
real $a$ is only the displayed integral on the same original measure.
Evenness of $\sigma$ and $t=y^2$ give, including both real half-lines,
\[
 \mu_a=\int_0^\infty t^{a-1/2}\sigma(\sqrt t)\,dt,
 \quad
 H_s^{(a)}=\frac1{s!}\int_{(0,\infty)^s}
       \prod_{i<j}(t_i-t_j)^2
       \prod_{i=1}^s t_i^{a-1/2}\sigma(\sqrt{t_i})\,dt_i.
 \tag{QGT3}
\]
To verify the second identity, expand the two determinants
$\det[t_i^{j-1}]_{i,j=1}^s$ as finite permutation sums, multiply them,
and integrate each monomial against the displayed positive product
measure. Every integral is finite: the density is bounded at zero,
$a>-1/2$, and the global Gamma upper bound below gives exponential
decay in $\sqrt t$ at infinity. Relabeling the integration variables
in each term gives $s!$ copies of the defining moment determinant.
The Vandermonde identity follows by its alternating zeros and its
leading coefficient one. This proves QGT3 with its literal factorial.
It is strictly positive because its integrand is positive on each
open cell of distinct positive coordinates.

\paragraph{The actual Gamma density, including its prefactor.}
Write $h=\pi/2$, $D=1/16+\sqrt3/18$, and set
\[
 C_+=\max\{\sigma(0)e^h,\sigma(1)e^{h+D}\},
 \qquad r(y)=y^{1/2}e^{hy}\sigma(y)\quad(y>0).
 \tag{QGT4}
\]
The accompanying proof GPR, obtained by integrating the proved original
logarithmic derivative GDR3, establishes on the full stated domains
\[
 \sigma(y)\ge\sigma(1)(1+y)^{-1/2}e^{-hy}\quad(y\ge0),
 \qquad \sigma(y)\le C_+y^{-1/2}e^{-hy}\quad(y>0).
 \tag{QGT5}
\]
Here is also the short full deduction needed for this use. GDR3 says
$L(y)=h+1/(2y)+\epsilon(y)/y$ and
$|\epsilon(y)|\le(1/8+\sqrt3/9)y^{-2}$. Integrating from $1$ to
$y\ge1$ shows
\[
 \log r(y)=\log\sigma(1)+h-\int_1^y\epsilon(u)\,du/u,
 \qquad \left|\int_1^y\epsilon(u)\,du/u\right|\le D.
\]
Since $h>D$, this gives $r(y)\ge\sigma(1)$ for $y\ge1$ and the
upper in QGT4 there. GDR7 implies $L(y)>0$ for $y>0$, so for
$0\le y\le1$ one has $\sigma(1)\le\sigma(y)\le\sigma(0)$.
On that interval $(1+y)^{-1/2}e^{-hy}\le1$ and
$y^{-1/2}e^{-hy}\ge e^{-h}$ when $y>0$. These facts prove both
inequalities in QGT5 and keep the density's finite value at zero.
The elementary inequality $h>D$ follows, for example, from
$\pi>3$ and $\sqrt3<2$.

Use exactly the dilation $t_i=q^2x_i$ in QGT3 and put
\[
 \alpha=\frac{a-3/4}{s},\qquad \beta=\frac{\pi q}{2s},
 \quad
 \mathcal J_{s,q}(\alpha,\beta)=\frac1{s!}
 \int_{(0,\infty)^s}\prod_{i<j}(x_i-x_j)^2
 \prod_i e^{-s(\beta\sqrt{x_i}-\alpha\log x_i)}
                     r(q\sqrt{x_i})\,dx_i.
 \tag{QGT6}
\]
Then the exact original determinant is
\[
 \boxed{H_s^{(a)}
   =q^{,2as+2s(s-1)+s/2}\mathcal J_{s,q}(\alpha,\beta).}
 \tag{QGT7}
\]
Indeed the Vandermonde contributes $q^{2s(s-1)}$, and each density
factor transforms as
\[
 (q^2x)^{a-1/2}\sigma(q\sqrt x)q^2dx
 =q^{2a+1/2}x^{a-3/4}e^{-hq\sqrt x}r(q\sqrt x)dx.
\]
This accounts separately for the original Jacobian, the quarter shift,
and the factor $q^{s/2}$. None of them is incorporated in the
equilibrium free energy.

\paragraph{A uniform finite partition theorem for this exact source.}
Let $\mathcal K=[3/2,5/2]\times[\pi/2,3\pi/2]$. Use the uniquely
proved equilibrium $\mu_{\alpha,\beta}$ of EIQ on the positive
half-line, its actual endpoints $a(\alpha,\beta),b(\alpha,\beta)$,
and its energy
\[
 F(\alpha,\beta)=\iint\log|x-y|\,d\mu_{\alpha,\beta}(x)
                    d\mu_{\alpha,\beta}(y)
       -\int(\beta\sqrt x-\alpha\log x)\,d\mu_{\alpha,\beta}(x).
 \tag{QGT8}
\]
In this paragraph only, the endpoint $a(\alpha,\beta)$ carries
arguments to distinguish it from the moment exponent $a$.
Let $a_*,B_*,M_*,C_0$ be the finite constants explicitly defined and
proved in QLG8--14 and QLG30. Set
\[
 L_0=\log(B_*M_*),\qquad
 C_* =\sigma(1)
          \left(\frac{\sqrt{a_*}}{1+\sqrt{a_*}}\right)^{1/2}>0,
\]
\[
 C=\max\{0,\ C_0+\log C_+,\ L_0-\log C_*\}.
 \tag{QGT9}
\]
For every $q\ge1$, integer $s\ge1$ and $(\alpha,\beta)\in\mathcal K$,
\[
 \boxed{\left|\log\mathcal J_{s,q}(\alpha,\beta)
                         -s^2F(\alpha,\beta)\right|\le Cs.}
 \tag{QGT10}
\]
For the upper estimate QGT5 gives $r(q\sqrt x)\le C_+$ everywhere
on the original positive half-line. QLG31 for the Lebesgue reference
measure therefore proves
$\log\mathcal J\le s^2F+(C_0+\log C_+)s$.
For the lower estimate retain the precise quantile cells QLG24--28
for this same equilibrium and its same $s$. Every coordinate in
these cells lies in $[a_*,B_*]$, where QGT5 gives
\[
 r(q\sqrt x)\ge\sigma(1)
      \left(\frac{q\sqrt x}{1+q\sqrt x}\right)^{1/2}\ge C_*.
\]
The last inequality uses $q\ge1$, $x\ge a_*$ and monotonicity of
$u/(1+u)$. QLG25--28 bound the integral over precisely these cells,
with their $s!$ permutations cancelling the original Heine factor.
Multiplying its integrand by the lower factor $C_*^s$ yields
$\log\mathcal J\ge s^2F-(L_0-\log C_*)s$.
This proves QGT10. A global lower bound by $x^{-3/4}$ at zero was
not used; the original density and its exact positive-domain factor
$r(q\sqrt x)$ remain in QGT6.

Combining QGT7 and QGT10 gives the exact finite decomposition
\[
 \boxed{\log H_s^{(a)}=
 [2as+2s(s-1)+s/2]\log q
 +s^2F\left(\frac{a-3/4}{s},\frac{\pi q}{2s}\right)+E_{s,q}(a),
 \qquad |E_{s,q}(a)|\le Cs.}
 \tag{QGT11}
\]
The bound holds whenever the displayed parameter pair belongs to
$\mathcal K$. The quantity $E_{s,q}(a)$ denotes the exact difference
in this equality; no derivative estimate for it has been assumed.

\paragraph{A quantitative moving secant, proved from convexity.}
Write $\mathcal L(\alpha,\beta)=\partial_\alpha F(\alpha,\beta)$,
the smooth original logarithmic moment proved in EIQ26 and EIQ29--30.
The fixed-interval integral EIQ29 and its proved differentiated
majorants give a finite constant
\[
 M=\max_{(\alpha,\beta)\in\mathcal K}
                   |\partial_\alpha\mathcal L(\alpha,\beta)|,
 \qquad A=M/2+2C.
 \tag{QGT12}
\]
For the two original ranks $s=n,n+1$, every real initial exponent
$a\in[q,q+l+1]$, and the original positive integer shift $l$, one has
\[
 \boxed{\left|\log\frac{H_s^{(a+l)}}{H_s^{(a)}}-2sl\log q
 -s^2\left[F\left(\frac{a+l-3/4}{s},\frac{\pi q}{2s}\right)
           -F\left(\frac{a-3/4}{s},\frac{\pi q}{2s}\right)\right]
                         \right|\le A l\sqrt s.}
 \tag{QGT13}
\]
All parameter neighbourhoods used to prove this bound lie in the
stated rectangle. Indeed put $u=(a-3/4)/s$, $t=l/s$ and
$\eta=s^{-1/2}<1/4$. For either rank,
\[
 u\ge2-\frac{11}{4s}\ge2-\frac{11}{72},\qquad
 u+t\le2+\frac{2l+1/4}{n}\le2+\frac18.
\]
Here $n\ge18l$ and $n\ge18$ give
$(2l+1/4)/n\le1/9+1/72=1/8$. Thus
$[u-\eta,u+t+\eta]\subset(3/2,5/2)$.
The second coordinate is $\pi$ for $s=n$ and $\pi n/(n+1)$ for
$s=n+1$, both lying in $[\pi/2,3\pi/2]$.

For fixed $q,s,\beta$, define
$g(v)=s^{-2}\log\mathcal J_{s,q}(v,\beta)$ on this neighbourhood.
This function is twice continuously differentiable and convex.
To justify differentiation under its integral, on every compact
subinterval of $v>0$ the exponent at zero is a positive power of
$x$, and the Gamma upper in QGT5 bounds the tails by a polynomial
times $\exp(-s\beta\sqrt x)$. The finite Vandermonde polynomial,
multiplied by any of the finitely many $\log x_i$ or their products
arising in the first two derivatives, has an integrable majorant.
Division by the strictly positive integral is legitimate. Its second
logarithmic derivative is $s^2$ times the variance of
$\sum_i\log x_i$ under its exact probability density and is
nonnegative. This proves convexity without a property of the remainder.

For each $v\in[u,u+t]$, convexity gives
\[
 \frac{g(v)-g(v-\eta)}{\eta}\le g'(v)
              \le\frac{g(v+\eta)-g(v)}{\eta}.
\]
QGT10 implies $|g-F|\le C/s$. The integral formula for a difference
quotient of $F$ and the bound in QGT12 give, for both signs,
\[
 |g'(v)-\mathcal L(v,\beta)|
          \le\frac{M\eta}{2}+\frac{2C}{s\eta}.
 \tag{QGT14}
\]
In detail, the upper quotient of $F$ is
$\eta^{-1}\int_v^{v+\eta}\mathcal L(w,\beta)dw$ and differs
from $\mathcal L(v,\beta)$ by at most
$\eta^{-1}\int_0^\eta Mr\,dr=M\eta/2$; the lower quotient
has the identical bound. Integrating QGT14 over $[u,u+t]$ and
multiplying by $s^2$ gives $Al\sqrt s$ because
$\eta=s^{-1/2}$ and $t=l/s$. In QGT7 the terms
$2s(s-1)\log q$ and $(s/2)\log q$ cancel exactly between the
two values of $a$, while $2as\log q$ contributes $2sl\log q$.
This proves QGT13, retaining the rank $n+1$ and the source prefactor.

\paragraph{All four original high blocks in one exact centre.}
Define the same positive high product as in LET and LER,
\[
 Z_a=H_n^{(a)}H_{n+1}^{(a)}(H_n^{(a+1)})^2.
\]
For the original pair of endpoints $a=q,q+l$ put
\[
\begin{split}
 \mathcal C_{q,l}={}&n^2\left[
 F\left(\frac{q+l-3/4}{n},\pi\right)
 -F\left(\frac{q-3/4}{n},\pi\right)\right]\\
 &+(n+1)^2\left[
 F\left(\frac{q+l-3/4}{n+1},\frac{\pi n}{n+1}\right)
 -F\left(\frac{q-3/4}{n+1},\frac{\pi n}{n+1}\right)\right]\\
 &+2n^2\left[
 F\left(\frac{q+l+1/4}{n},\pi\right)
 -F\left(\frac{q+1/4}{n},\pi\right)\right],
\end{split}
 \tag{QGT15}
\]
\[
 \mathcal U_{q,l}=Al(3\sqrt n+\sqrt{n+1}).
 \tag{QGT16}
\]
Adding QGT13 with the displayed coefficients $1,1,2$ yields the
finite exact decomposition
\[
 \boxed{\log\frac{Z_{q+l}}{Z_q}
       =(4q+2)l\log q+\mathcal C_{q,l}+\xi_{q,l},
              \qquad |\xi_{q,l}|\le\mathcal U_{q,l}.}
 \tag{QGT17}
\]
The coefficient of the logarithm is exactly
$2l[n+(n+1)+2n]=2l(2q+1)$. The two odd blocks, the different
high ranks, and the original quarter shifts are all explicit in QGT15.

\paragraph{Finite return with the low endpoint and every product correction.}
Use the original $P_k$ of WGP1 and the original scalar moment of LER1,
so that
\[
 W_k=P_k+\log(Z_{q+l}/Z_q)-\log(\mu_{q+l}/\mu_q).
 \tag{QGT18}
\]
Put $\widetilde P_{q,l}=P_k+4lq\log q$, the exact finite sum WGP3,
and retain $F_{q,l},K_{q,l},T_{q,l},E_{q,l}$ with exactly the
definitions LER13 and LER17. Define the following centre in terms of
the already proved original objects:
\[
 \mathcal V_{q,l}=\widetilde P_{q,l}+\mathcal C_{q,l}
                         +2l\log q-F_{q,l}+K_{q,l}.
 \tag{QGT19}
\]
QGT17 and the signed remainder LER18 give
\[
 \boxed{\mathcal V_{q,l}-\mathcal U_{q,l}-E_{q,l}
       \le W_k\le
       \mathcal V_{q,l}+\mathcal U_{q,l}+T_{q,l}+E_{q,l}.}
 \tag{QGT20}
\]
To check the signs, LER18 is
$\log(\mu_{q+l}/\mu_q)=F_{q,l}-K_{q,l}+r_{q,l}+\varepsilon_{q,l}$
with $-T_{q,l}\le r_{q,l}\le0$ and
$|\varepsilon_{q,l}|\le E_{q,l}$. Substitution in QGT18 gives
$W_k=\mathcal V_{q,l}+\xi_{q,l}-r_{q,l}-\varepsilon_{q,l}$,
which proves QGT20. The original logarithmic powers cancel exactly:
$-4lq+(4q+2)l-2l=0$, since the logarithmic part of $F_{q,l}$ is
$2l\log(4q/\pi)$.

For an entirely integral-and-rational centre instead of the finite
source sum, set
\[
 p^{\rm c}_{q,l}=-4lq(2\log2-1)-4l^2\log2
       +\frac{8l^3+l}{6q}-\frac{l^4+l^2/4}{q^2},
\]
\[
 p^-_{q,l}=\frac l{6q}+\frac{8l^3+l}{12q^3},\qquad
 p^+_{q,l}=\frac{l^2}{3q^2}
       +\frac{8l^4+2l^2}{6q^4}
       +\frac{7(192l^5+80l^3-2l)}{1440q^3},
\]
\[
 \mathcal V^{\rm c}_{q,l}
     =p^{\rm c}_{q,l}+\mathcal C_{q,l}
                        +2l\log q-F_{q,l}+K_{q,l}.
 \tag{QGT21}
\]
WGP13 is the proved finite statement
$-p^-_{q,l}\le\widetilde P_{q,l}-p^{\rm c}_{q,l}\le p^+_{q,l}$.
Consequently the full original Gamma functional satisfies
\[
 \boxed{\mathcal V^{\rm c}_{q,l}-p^-_{q,l}-\mathcal U_{q,l}-E_{q,l}
       \le W_k\le
       \mathcal V^{\rm c}_{q,l}+p^+_{q,l}
                         +\mathcal U_{q,l}+T_{q,l}+E_{q,l}.}
 \tag{QGT22}
\]
In particular this bound uses the exact Gamma source and both original
low and high endpoint contributions; its finite constants were
obtained from their integrals and the proved equilibrium.

\paragraph{An explicit bound around the previously proved leading coefficient.}
Write $\mathcal L_0=\mathcal L(2,\pi)$ and
$\mathcal C_\Gamma=2\mathcal L_0-4(2\log2-1)$, exactly the positive
coefficient proved in GEL. The actual dilation proved in EIQ32 gives
$\mathcal L(v,\pi n/(n+1))=
\mathcal L(v,\pi)+2\log(1+1/n)$.
For $d\in\{-3/4,-11/4,1/4\}$, the fundamental theorem of calculus
and QGT12 give
\[
 \left|s^2\int_{2+d/s}^{2+(d+l)/s}
           \mathcal L(v,\pi)\,dv-sl\mathcal L_0\right|
 \le M\left(|d|l+\frac{l^2}{2}\right).
 \tag{QGT23}
\]
Indeed set $v=2+(d+r)/s$, $0\le r\le l$; the integrand's
difference from $\mathcal L_0$ is at most $M|d+r|/s$,
and $|d+r|\le |d|+r$. Apply this respectively at ranks $n,n+1,n$
with coefficients $1,1,2$ in QGT15. Since
$3/4+11/4+2(1/4)=4$, the result is
\[
 \left|\mathcal C_{q,l}-(2q+1)l\mathcal L_0
                      -2(n+1)l\log(1+1/n)\right|
                    \le M(2l^2+4l).
 \tag{QGT24}
\]
Every $v$ in these integrals is in the rectangle already verified above.
Define the fully explicit finite quantity
\[
\begin{split}
 \mathcal D_{q,l}={}&M(2l^2+4l)+l|\mathcal L_0|
        +2(n+1)l\log(1+1/n)+4l^2\log2\\
 &+\frac{8l^3+l}{6q}+\frac{l^4+l^2/4}{q^2}
        +\max\{p^-_{q,l},p^+_{q,l}\}
        +2l|\log(4/\pi)|+\frac{l^2}{q}\\
 &+\frac{l(16l^2-1)}{48q^2}
        +\frac{l^2(8l^2-1)}{48q^3}
        +K_{q,l}+T_{q,l}+E_{q,l}.
\end{split}
 \tag{QGT25}
\]
All displayed terms are nonnegative for $l\ge1$. Combining
QGT18, QGT17, WGP13, LER17--18 and QGT24, with the triangle inequality
applied only after their signs have been recorded, proves
\[
 \boxed{|W_k-lq\mathcal C_\Gamma|
                     \le\mathcal U_{q,l}+\mathcal D_{q,l}.}
 \tag{QGT26}
\]
The sign-sensitive sharper statement remains QGT22; QGT26 is its
explicit estimate about the original leading coefficient.

Along the actual packet, $\mathcal D_{q,l}=O(l^2+l)$ with a constant
independent of $m$. This follows term by term from
$q\ge(4l+2)^2$, the displayed rational quantities, LER20, and
$\log(1+1/n)\le1/n$ (integrate $1/(1+t)\le1$ on $[0,1/n]$).
Furthermore $\mathcal U_{q,l}\le4Al\sqrt{n+1}$.
Thus the quantitative theorem proves, for the same original integers,
\[
 \boxed{W_k=lq\mathcal C_\Gamma+O(l\sqrt q+l^2).}
 \tag{QGT27}
\]
The finite bounds QGT22 and QGT26 specify this error whenever an
order symbol is insufficient.
For each fixed $m\ge2$, the exact formula
$q=[1+(4l+1)(m-1)](4l+2)^2$ gives
$l/\sqrt q\to0$ and $l^2/q\to0$. Dividing QGT26 by $q$ therefore
proves the stronger original-return limit
\[
 \boxed{\frac{W_k-lq\mathcal C_\Gamma}{q}\longrightarrow0
                        \quad(l\to\infty,\ m\ge2\text{ fixed}).}
 \tag{QGT28}
\]
For $m=1$, the same exact calculation gives a finite bound on
$(W_k-lq\mathcal C_\Gamma)/q$, since $q=(4l+2)^2$ and
$l/\sqrt q\to1/4$. The bound does not identify the limit of that
quotient: the retained high-block uncertainty
$\mathcal U_{q,l}/q$ need not tend to zero in this case.
This identifies precisely what this finite estimate resolves at the
original $q$ scale, with every packet relation unchanged.
